% Options for packages loaded elsewhere
\PassOptionsToPackage{unicode}{hyperref}
\PassOptionsToPackage{hyphens}{url}
%
\documentclass[
]{article}
\usepackage{amsmath,amssymb}
\usepackage{iftex}
\ifPDFTeX
  \usepackage[T1]{fontenc}
  \usepackage[utf8]{inputenc}
  \usepackage{textcomp} % provide euro and other symbols
\else % if luatex or xetex
  \usepackage{unicode-math} % this also loads fontspec
  \defaultfontfeatures{Scale=MatchLowercase}
  \defaultfontfeatures[\rmfamily]{Ligatures=TeX,Scale=1}
\fi
\usepackage{lmodern}
\ifPDFTeX\else
  % xetex/luatex font selection
\fi
% Use upquote if available, for straight quotes in verbatim environments
\IfFileExists{upquote.sty}{\usepackage{upquote}}{}
\IfFileExists{microtype.sty}{% use microtype if available
  \usepackage[]{microtype}
  \UseMicrotypeSet[protrusion]{basicmath} % disable protrusion for tt fonts
}{}
\makeatletter
\@ifundefined{KOMAClassName}{% if non-KOMA class
  \IfFileExists{parskip.sty}{%
    \usepackage{parskip}
  }{% else
    \setlength{\parindent}{0pt}
    \setlength{\parskip}{6pt plus 2pt minus 1pt}}
}{% if KOMA class
  \KOMAoptions{parskip=half}}
\makeatother
\usepackage{xcolor}
\setlength{\emergencystretch}{3em} % prevent overfull lines
\providecommand{\tightlist}{%
  \setlength{\itemsep}{0pt}\setlength{\parskip}{0pt}}
\setcounter{secnumdepth}{-\maxdimen} % remove section numbering
\ifLuaTeX
  \usepackage{selnolig}  % disable illegal ligatures
\fi
\usepackage{bookmark}
\IfFileExists{xurl.sty}{\usepackage{xurl}}{} % add URL line breaks if available
\urlstyle{same}
\hypersetup{
  hidelinks,
  pdfcreator={LaTeX via pandoc}}

\author{}
\date{}

\begin{document}

\section{THE MODEL COULD DRAW YELMO. IT COULD NOT KEEP
HIM.}\label{the-model-could-draw-yelmo.-it-could-not-keep-him.}

\subsection{Scott Roberts, ``the same dude each time,'' and the minimum
conditions of generative
storytelling}\label{scott-roberts-the-same-dude-each-time-and-the-minimum-conditions-of-generative-storytelling}

\textbf{Watson Hartsoe · Working Paper · 17 August 2026}

\subsubsection{Abstract}\label{abstract}

The original Yelmo project should not be remembered as a failed attempt
to make a comic with Midjourney. Its proposal was more basic and more
consequential: could an image generator sustain a protagonist long
enough to tell a story? The surviving 2022 notes describe a small group
imagining ``a hero's journey of Midjourney,'' searching for ``someone
that's kind of repeatable,'' and settling on a deliberately wrong but
recognizable figure: yellow Elmo, quickly compressed into Yelmo. The
joke became a stress test. Across scenes, Yelmo's eyes, frill, body, and
relations drifted; adding other characters could cause them to alter one
another. The problem was not simply visual inconsistency. The system
could repeatedly generate descriptions of a type, but it did not carry
forward a stable referent --- \emph{this same character} --- from one
image to the next.

This paper argues that Yelmo exposed a minimum ontology of generative
storytelling before the field had good names for it. A story requires
identity that survives transformation, not mere repetition; multiple
actors whose boundaries survive interaction; and some mechanism that
remembers what the story has already committed to. In the Yelmo
experiments, much of that state appears to have lived outside the model,
in the human group that named variants, rejected mutations, remembered
prior versions, and treated a drifting family of images as one
protagonist. \textbf{The model could draw Yelmo. The community was
keeping Yelmo alive.}

The argument becomes more interesting beside Scott Roberts's documented
practice. Years before generative AI, Roberts made work about entering
the world of a cereal mascot, animating gamers' idealized alter egos,
crossing representation and physical space, and ``experiment{[}ing{]}
while making comics.'' He later published a chapter titled ``Surrealist
Storytelling: Artistic Experimentation and Spontaneity through
Improvisation,'' publicly described Midjourney as a medium for an
ongoing exquisite-corpse game in August 2022, and in 2023 taught an
experimental AI-art course whose students he called research
collaborators. None of this proves a direct lineage from those works to
Yelmo. It does make a stronger question possible: was Yelmo not an AI
novelty at all, but a new machine entering an older Roberts problem ---
how a character survives translation, improvisation, and transformation?

\textbf{Keywords:} generative AI; Midjourney; animation; character
consistency; visual storytelling; improvisation; subject identity; Scott
Roberts

\begin{center}\rule{0.5\linewidth}{0.5pt}\end{center}

\subsection{1. The joke was already a theory of
identity}\label{the-joke-was-already-a-theory-of-identity}

The original proposal, as it survives in the research notes, is almost
disarmingly small. A few people are trying to imagine ``a hero's journey
of Midjourney.'' They need a protagonist --- ``someone that's kind of
repeatable.'' A friend named Scott proposes, or helps develop, an absurd
test figure: yellow Elmo in a gritty farm-documentary register. The
reason is partly comic. Elmo is not yellow. ``Yellow Elmo'' rhymes its
way into a name. Soon the group has Yelmo hiking a mountain, sitting at
a fancy dinner, wandering into whatever scenario somebody can throw at
him.

That choice is better than it looks.

Yelmo begins as a contradiction that remains recognizable. He is wrong
before the experiment starts: Elmo recolored, category violated,
identity bent. Yet he is still recognizable enough for the joke to land.
The group has therefore chosen a subject that asks an unusually clean
question: \textbf{how much can a figure change before it stops being the
figure?}

That is not the usual benchmark for an image generator. Image generators
are usually flattered by single-frame evaluation. Is this image
beautiful? Does it match the prompt? Does the hand have five fingers? Is
it photorealistic? Yelmo refuses that generosity. A gorgeous Yelmo on a
mountain is not enough. The next image has to be \emph{him again}.

The surviving notes state the failure with more precision than most
technical vocabulary of the period: ``it's not really the same dude each
time.'' Sometimes the eyes change. Sometimes a red frill appears or
disappears. The model seems to merge different yellow-Elmo-like
possibilities. The image is good; the character is unstable.

That distinction is the beginning of the paper. \textbf{A model can
succeed at depiction while failing at reference.} It can produce an
image satisfying the description \emph{yellow Elmo in an Italian
restaurant} without possessing anything equivalent to the narrative
commitment \emph{the individual we already met on the mountain is now at
this table}.

The difference sounds philosophical because it is. A story does not
merely need recurring resemblance. It needs a way to say \emph{this one
again}.

\subsection{2. The original proposal was not ``make images.'' It was
``make a protagonist
survive.''}\label{the-original-proposal-was-not-make-images.-it-was-make-a-protagonist-survive.}

The Joseph Campbell frame matters more than the notes give it credit
for. A hero's journey is not a slideshow of unrelated illustrations. Its
basic wager is continuity through transformation. The hero leaves home,
crosses thresholds, meets allies and enemies, is altered by ordeals, and
returns changed. If the hero were visually frozen, nothing would have
happened. If the hero became a different person every frame, there would
be no one to whom the happening belonged.

So the original Yelmo proposal accidentally chose a much harder problem
than ``character consistency.'' \textbf{It asked for identity-preserving
transformation.}

That phrase matters. Consistency is often imagined as sameness: preserve
the face, the shirt, the fur, the proportions. Narrative asks for
something subtler. Change the pose. Change the lighting. Change the
costume. Injure him. Age him. Frighten him. Put him underwater. Let him
lose the coat that seemed visually definitive in the first image. Make
him return from the underworld with something about him genuinely
altered. And still: same dude.

This is why Yelmo is a better storytelling test than a static character
sheet. The project forces a system to decide which differences are
allowed to vary and which differences would destroy identity. Later
subject-driven generation research would formalize versions of exactly
this problem. Textual Inversion asked how a unique user-provided concept
could be represented and recomposed into new scenes. DisenBooth
explicitly separated identity-relevant information from pose,
background, and other identity-irrelevant information. ConsiStory later
pursued consistency across diverse prompts without requiring a
permanently trained subject representation. These systems disagree about
where persistence should live, but they share the problem Yelmo
encountered experientially: a subject must survive variation without
becoming trapped by its first depiction (Gal et al.~2022; Chen et
al.~2023; Tewel et al.~2024).

The strange thing is that there is no neutral answer to what counts as
identity-relevant. If every surviving Yelmo image has the same red
frill, the frill begins to feel essential. If one image loses it and the
group still says ``that's Yelmo,'' the group has demoted it to costume.
The reference set does not merely reveal identity; it helps legislate
identity.

The original project therefore contains a question that should be put
directly to Scott as an animator: \textbf{when you said ``the same dude
each time,'' what did you actually mean by same?} Which properties did
your eye forgive? Which mutations killed the character? Did you think in
terms of model-sheet invariants --- silhouette, facial spacing,
proportion, costume, color --- or was recognition looser and more
performative? Was a bad drawing still Yelmo if the joke, pose, or
attitude survived?

Those answers would matter more than another general interview about AI
creativity. They would recover the human criterion against which the
machine was being tested.

\subsection{3. The community was the
rig}\label{the-community-was-the-rig}

Here is the braver claim.

In the workflow described by the notes, \textbf{Yelmo did persist. He
just did not persist inside Midjourney.}

He persisted in the channel.

The people remembered what he had looked like. They named him. They
compared new generations with prior generations. They laughed when the
model produced an especially strange deviation. They rejected some
mutations and accepted others. They proposed new situations for the same
supposed individual. The continuity relation --- \emph{this image
belongs to the Yelmo lineage} --- was being maintained socially.

An animator will immediately see the inversion. A conventional character
pipeline externalizes continuity into model sheets, turnarounds,
palettes, rigs, reference drawings, naming conventions, asset libraries,
and human review. Those devices make the character portable across shots
and across hands. Yelmo had almost none of that inside the generative
system. The group became the continuity apparatus.

\textbf{The community was the rig.}

This is not a metaphor for collaboration. It is a description of where
state lived. The generator could be asked again for ``yellow Elmo,'' but
the social group carried the history of what prior outputs had already
established. Humans supplied the memory against which the next image
became a continuation, mutation, or failure.

This makes the famous instability of early image generators newly
interesting. Their lack of persistence did not simply prevent narrative.
It recruited humans into maintaining the missing persistence themselves.
The limitation helped produce the social game. Every broken Yelmo
created an occasion to say: Is that him? Is this better? What if we try
this? What if he goes here? The defect became a generator of
conversation.

That may explain why the project could be simultaneously a failure as a
production tool and a success as a social practice. The notes call
Midjourney a ``multiplayer imagination game.'' Participants riff on one
another, propose increasingly ridiculous scenarios, and watch the model
make moves they would not have proposed themselves. The system is bad at
maintaining a production asset and excellent at producing provocations.

The project's own language catches the split: \textbf{game, not tool}.
Not because nothing useful is happening, but because the labor of
coherence still belongs to the people. A tool can be inserted into a
pipeline because its outputs can be made answerable to prior decisions.
A game can be enjoyable precisely because prior decisions remain
vulnerable to surprise.

Yelmo sits at that threshold.

\subsection{4. The second character reveals the missing
world}\label{the-second-character-reveals-the-missing-world}

The strongest technical clue in the surviving notes is not the changing
frill. It is the failure that appears when another character enters.

``We couldn't add other characters.'' Yelmo would begin altering them;
the characters contaminated one another.

That is a different problem from one-character persistence. It suggests
that the model was not merely failing to remember Yelmo. It was failing
to maintain \textbf{independent entity boundaries} inside a shared
scene.

A story requires relations. Yelmo needs an ally, antagonist, lover,
stranger, mentor, crowd. But a relation is possible only if the relata
survive the relation. If introducing a second figure causes attributes
to leak until the figures converge, then the system does not yet have
two stable actors; it has one image field negotiating a compromise.

Later research on multi-subject generation gives this failure
surprisingly literal names. IR-Diffusion identifies interference among
subjects and proposes ``Isolation Attention'' to prevent one subject
from referencing another in ways that produce subject fusion. StoryBooth
describes ``self-attention leakage'' in which one character begins to
look like another and frames the problem specifically as an obstacle to
visual storytelling. These papers are not evidence that the Yelmo group
anticipated their mechanisms, and there is no influence claim here. They
are valuable because they let us rename the 2022 practical complaint (He
et al.~2024; Singh et al.~2025).

\textbf{Yelmo was discovering that a story needs object boundaries.}

A beautiful image can blur identities and remain beautiful. A story
cannot. If Yelmo shakes hands with a pirate and comes away half-pirate,
the failure is not cosmetic. The world has lost the ability to say who
did what to whom.

This is where the original project becomes more ambitious than most
discussions of prompting. The minimum generative story world requires at
least three operations:

\begin{enumerate}
\def\labelenumi{\arabic{enumi}.}
\tightlist
\item
  preserve an entity across transformations;
\item
  preserve several entities simultaneously;
\item
  allow relations among those entities without fusing their identities.
\end{enumerate}

The Yelmo group found all three by trying to make a ridiculous Muppet go
on an adventure.

\subsection{5. Scott Roberts makes the story harder, not
easier}\label{scott-roberts-makes-the-story-harder-not-easier}

There is still a provenance problem that should not be hidden because
Scott may be the person who can solve it.

The Yelmo transcript identifies a ``Scott,'' describes him as an
animator and the more professionally experienced artist in the group,
and places him inside the project. The researcher supplying this archive
identifies that Scott as DePaul animator and professor Scott Roberts.
Public evidence independently places Roberts deep in Midjourney
experimentation in the same period: on August 31, 2022, he described
work made in Midjourney that included results from an ``on-going
exquisite corpse game''; a collaborator had publicly called another
August exchange an ``exquisite corpse collaboration'' with Roberts
through Midjourney. But I have not found a public Roberts source that
says ``Yelmo.'' That bridge should remain visible until Roberts confirms
it or an original record does.

The interesting part is that Roberts's documented pre-AI work makes
Yelmo more, not less, plausible as serious artistic research.

In 2002, Roberts made \emph{Lucky}, a projection sculpture that he
describes as allowing him to ``enter the world of the cereal mascot,'' a
childhood dream. The same year, \emph{My Favorite Character Is a Wizard}
drew from interviews with role-playing gamers about their favorite
characters; Roberts says the animation brings those ``idealized alter
egos to life'' while revealing the mannerisms and idiosyncrasies of
their creators. His artworks repeatedly cross representational systems:
a cartoon cat projected into architecture, a low-resolution digital
model texture-mapped with high-resolution photographs, objects conceived
as Boolean intersections, drawings used to imitate the slow movement of
a motion-picture camera. His comics page opens with the plain
declaration: ``I experiment while making comics'' (Roberts, ``Comics'').

Then there is the later chapter title: ``Surrealist Storytelling:
Artistic Experimentation and Spontaneity through Improvisation.'' We do
not yet have the chapter body in this archive, so its argument cannot be
smuggled in from the title. But the conjunction is difficult to ignore:
storytelling, experiment, spontaneity, improvisation. Those are also the
terms needed to describe the Yelmo practice at its best.

This does not prove Yelmo descended from \emph{Lucky}, the Wizard
animation, or a Surrealist method. Chronology cannot do that work. But
it gives us a better historical question than ``When did Scott start
using AI?''

The better question is: \textbf{what old problems did Midjourney
suddenly make cheap enough, fast enough, and social enough for Roberts
to attack in a new way?}

The mascot matters. The alter ego matters. The representation crossing
into another world matters. The willingness to improvise matters. If
Yelmo is Roberts's project, then the most interesting reading may be
that generative AI did not create a new artistic concern. It gave an old
concern a machine partner that could mutate a character faster than a
human group could stabilize him.

\subsection{6. The hero's journey failed into a better research
project}\label{the-heros-journey-failed-into-a-better-research-project}

The notes suggest that the initial Joseph Campbell project did not
harden into a finished narrative. It ``morphed into just developing that
character.'' This is usually the point where a project history becomes
apologetic: the story was never completed; the technology was not ready;
the group got distracted.

That is the wrong reading.

The story proposal failed into the experiment it actually needed to
become.

Before asking whether Midjourney could tell twelve stages of a myth, the
group had to establish whether the same hero could survive stage one and
stage two. Yelmo reduced ``Can AI tell a story?'' to a prior question:
\textbf{Can there be a recurring somebody for the story to happen to?}

That is methodological progress. It takes a vague cultural question and
turns it into a testable production problem.

The project then found the next prior question: can the hero encounter
somebody else without identities bleeding together? And another: which
visual features are constitutive of identity, and which can change? And
another: if the machine is most enjoyable when it surprises the group,
how much control can be added before the improvisational value
disappears?

This last tension is especially important for Roberts because the notes
preserve both sides. On one side is the animator's demand for
intentionality: enough control to tell a story, define a character, and
carry work into the world. On the other side is the pleasure of ``being
led by MJ,'' of not trying to control everything, of receiving a visual
move no artist in the group would have proposed.

That is not a temporary contradiction to be solved by better software.
It is a durable artistic problem.

\textbf{The more perfectly the system obeys, the less it can surprise.
The more radically it surprises, the less it can sustain a world.}

Yelmo is valuable because it places those demands in the same body. Keep
him. Change him. Surprise me. Do not lose him.

That is a better research program than ``AI-assisted animation.''

\subsection{7. Midjourney eventually named the
requirement}\label{midjourney-eventually-named-the-requirement}

There is a clean historical afterimage.

Midjourney's own documentation describes a Character Reference as a way
to recreate a specific character in multiple images and maintain
consistency across scenes. The company announced Character References on
March 12, 2024 (Midjourney 2024; Midjourney, ``Character Reference'').
In later versions, Omni Reference similarly lets users carry characters
and objects into new creations. Midjourney's documentation now presents
as a product feature almost exactly the capability the Yelmo group was
missing in 2022.

No causal claim follows. There is no evidence here that Yelmo influenced
Midjourney's roadmap. The important point is convergence. An early
practitioner group independently formulated a limitation that later
became legible enough to receive a dedicated interface and product
vocabulary.

The original notes even contain the product argument before the product
existed: without character consistency, the images remain ``a beautiful
garden of images'' trapped as one-offs; with a definable character that
can be placed in new situations, comics, animation, poster series, and
sustained storytelling open up.

That is more than user frustration. It is a requirements document spoken
in the vernacular of play.

The braver historical claim is therefore this: \textbf{Yelmo was an
informal benchmark before the product had a benchmark for the thing that
mattered.}

Not image quality. Not beauty. Not prompt eloquence.

Can the character come back?

\subsection{8. The experiment should be rerun, not
memorialized}\label{the-experiment-should-be-rerun-not-memorialized}

If the original Yelmo materials survive, the project is unusually well
suited to reconstruction.

The goal should not be to make a nostalgic gallery. It should be to
rerun the same problem across the history of the medium.

Recover the original prompts, outputs, timestamps, remix chains,
private-server messages, and the October 12, 2022 Office Hours material
noted in the archive. Recover the earliest image the group agreed to
call Yelmo. Recover the first moment a participant said some variant was
\emph{not} him. Recover the failed multi-character images. Recover the
abandoned Campbell sequence, if it ever existed as a list of stages or
prompts.

Then build the benchmark the group was accidentally inventing:

\textbf{Persistence.} Put Yelmo through twelve sharply different scenes.
Ask whether independent raters believe the same character survives.

\textbf{Transformation.} Deliberately alter pose, costume, affect, age,
lighting, medium, and environment. Determine which transformations
preserve identity and which destroy it.

\textbf{Re-entry.} Remove Yelmo for several scenes, then bring him back.
Does the system recover the same individual or merely another member of
the type?

\textbf{Multiplicity.} Add one, two, then five recurring characters.
Measure when identities begin to contaminate one another.

\textbf{Interaction.} Make characters touch, exchange objects, occlude
each other, change clothes, or share props. This is where mere visual
separation becomes world coherence.

\textbf{Improvisation.} Run two conditions: one optimized for maximum
identity control, another allowing the generator much more freedom.
Compare not just fidelity but the number of genuinely useful surprises
introduced by the machine.

This would turn Yelmo into something rare: a longitudinal artistic
benchmark whose question survives radical changes in model architecture.

The point would not be to prove that 2026 tools are better than 2022
tools. Of course they are better at many of these operations. The point
would be to preserve the original human criterion and watch the medium
slowly become able to satisfy it.

\subsection{9. Scott, this is the part only you can
break}\label{scott-this-is-the-part-only-you-can-break}

If the Scott in the Yelmo record is Scott Roberts, then the next step is
not another literature review. It is to let you damage this argument.

The paper needs answers to questions only the original participants can
supply:

\textbf{What was the proposal before Yelmo existed?} Was there actually
a planned Campbell sequence, or is ``hero's journey'' a retrospective
shorthand for a looser desire to make a story?

\textbf{Who invented Yelmo?} Who first typed the yellow-Elmo prompt, who
named him, and was the fact that Elmo ``should'' be red part of the
conceptual attraction or simply a joke that worked?

\textbf{What did repeatable mean to you as an animator?} What were the
non-negotiable features? Which variations did you accept as performance,
costume, or bad drawing, and which made you say: no, that is another
character?

\textbf{Did you ever make an actual model sheet?} Did the group try to
formalize Yelmo outside the prompt --- reference boards, prompt
templates, seeds, image prompts, Photoshop corrections, naming rules,
canonical images?

\textbf{When did the hero's journey become the Yelmo laboratory?} Was
there a particular failure that killed the narrative plan, or did the
social pleasure of putting him in random situations simply become the
better project?

\textbf{What happened when other characters entered?} Did you understand
the failure as prompt competition, model confusion, compositing failure,
or something else? Do the images survive?

\textbf{How much did you want to lose control?} The notes contain both
the animator's frustration with insufficient intentionality and the
pleasure of letting Midjourney lead. Were those opposite camps in the
group, or two demands you personally wanted the tool to satisfy at once?

\textbf{Was exquisite corpse already your method?} Your August 2022
posts document Midjourney exquisite-corpse play, and your later
published chapter names Surrealist storytelling and improvisation. Did
that tradition consciously shape Yelmo, or is the resemblance something
this paper is imposing after the fact?

\textbf{Did Yelmo enter the 2023 DePaul course?} Your public description
of that course calls the students research collaborators exploring
bleeding-edge tools. Did Yelmo, character consistency, collaborative
prompting, or the game/tool distinction become examples or assignments?

\textbf{What do you still have?} The most valuable source may not be an
interview answer at all. It may be a folder of ugly images, a forgotten
Discord export, a screenshot of a prompt chain, or a note naming the
twelve stages.

If those materials contradict the paper, good. The point is not to turn
Yelmo into a myth. It is to recover what the experiment actually
discovered.

\subsection{Conclusion: the character existed in the channel before he
existed in the
model}\label{conclusion-the-character-existed-in-the-channel-before-he-existed-in-the-model}

The original Yelmo project asked a question more durable than the
technology that failed it.

Can a generative image system sustain a protagonist through change?

In 2022, the answer in this small experiment appears to have been: not
by itself. Midjourney could produce remarkable individual Yelmo images.
It could vary him endlessly. It could surprise the group. But it could
not reliably honor the accumulating commitments that made one generated
image the continuation of another. It could render a character
description without consistently maintaining a character history.

The humans closed the gap. They remembered. They selected. They named.
They laughed at mutations and carried accepted ones forward. They
supplied the distinction between variation and replacement. They held
the entity boundaries that broke when other characters entered. They
supplied enough continuity for Yelmo to become socially real even while
he remained computationally unstable.

That is why the original failure matters.

\textbf{The model could draw Yelmo. It could not keep him.}

And for a while, the people did.

If Roberts's older work belongs in this genealogy, the final irony is
even better. An artist who had already entered the world of a cereal
mascot, animated other people's idealized alter egos, crossed cartoons
into physical space, experimented through comics, and written about
improvisational storytelling encountered a machine that could summon
characters almost instantly but could not promise that the character
would still be there in the next scene.

The brave version of the history is not that generative AI suddenly
empowered an animator.

It is that an animator and his friends used a stupid joke to discover
what the machine still lacked before the field had agreed on the name of
the problem.

Yelmo was the joke.

``The same dude each time'' was the research question.

And the unfinished hero's journey may be the most important artifact
precisely because it never got past the condition a story needs before
the adventure can begin: \textbf{someone has to be able to leave one
image and arrive in the next.}

\begin{center}\rule{0.5\linewidth}{0.5pt}\end{center}

\subsection{Source note}\label{source-note}

The core Yelmo evidence is a researcher-supplied 2022
Midjourney-community transcript/notes file. It describes a ``Scott'' who
is an animator and participant in Yelmo but does not, in the surviving
excerpt, provide a surname. The researcher identifies this person as
Scott Roberts. Public sources independently document Roberts's
Midjourney practice in August 2022, including exquisite-corpse
collaboration, but no public source located for this revision explicitly
names Yelmo. The identification therefore remains a direct question for
Roberts rather than a concealed assumption.

\subsection{References}\label{references}

Chen, Hong, Yipeng Zhang, Simin Wu, Xin Wang, Xuguang Duan, Yuwei Zhou,
and Wenwu Zhu. 2023. ``DisenBooth: Identity-Preserving Disentangled
Tuning for Subject-Driven Text-to-Image Generation.'' arXiv:2305.03374.
https://arxiv.org/abs/2305.03374.

DePaul University. ``Scott Roberts.'' Faculty profile.
https://www.depaul.edu/faculty/scott-roberts.

DePaul University Center for Teaching and Learning. 2023. ``28th Annual
Teaching and Learning Conference.'' May 5, 2023.
https://resources.depaul.edu/teaching-commons/events/teaching-learning-conference/Pages/28th-Annual-Teaching-and-Learning.aspx.

Gal, Rinon, Yuval Alaluf, Yuval Atzmon, Or Patashnik, Amit H. Bermano,
Gal Chechik, and Daniel Cohen-Or. 2022. ``An Image is Worth One Word:
Personalizing Text-to-Image Generation using Textual Inversion.''
arXiv:2208.01618. https://arxiv.org/abs/2208.01618.

He, Huiguo, Qiuyue Wang, Yuan Zhou, Yuxuan Cai, Hongyang Chao, Jian Yin,
and Huan Yang. 2024. ``Improving Multi-Subject Consistency in
Open-Domain Image Generation with Isolation and Reposition Attention.''
arXiv:2411.19261. https://arxiv.org/abs/2411.19261.

Midjourney. 2024. ``Introducing Character References.'' March 12, 2024.
https://updates.midjourney.com/character-refs/.

Midjourney. ``Character Reference.'' Documentation.
https://docs.midjourney.com/hc/en-us/articles/32162917505293-Character-Reference.

Midjourney. ``Legacy Features.'' Documentation.
https://docs.midjourney.com/hc/en-us/articles/33329788681101-Legacy-Features.

Research archive. 2022. ``Yelmo / Midjourney community interview and
field notes.'' Researcher-supplied transcript/notes; unpublished.

Roberts, Scott. 2002. ``Lucky.'' Ubutopia.
https://www.ubutopia.com/lucky.htm.

Roberts, Scott. 2002. ``My Favorite Character Is a Wizard.'' Ubutopia.
https://ubutopia.com/wizard.htm.

Roberts, Scott. ``Comics.'' Ubutopia.
https://www.ubutopia.com/comics.html.

Roberts, Scott. 2022. Instagram post describing ``an on-going exquisite
corpse game'' made with Midjourney. August 31, 2022.
https://www.instagram.com/p/Ch6jYdhOw0j/.

Roberts, Scott. ``Surrealist Storytelling: Artistic Experimentation and
Spontaneity through Improvisation.'' In \emph{The Rose Metal Press Field
Guide to Graphic Literature}, edited by Kelcey Ervick and Tom Hart. Rose
Metal Press. Table of contents available at
https://rosemetalpress.com/wp-content/uploads/2023/02/FGtoGraphicLiteratureTableofContents.pdf.

Roberts, Scott. 2023. ``The Cyborg Muse: The Potential of AI Tools to
Enhance and Accelerate Creativity.'' LinkedIn / DePaul Center for
Teaching and Learning.
https://www.linkedin.com/posts/scott-roberts-3a63523\_28th-annual-teaching-and-learning-conference-activity-7059776060971769856-V8fR.

Singh, Jaskirat, Junshen Kevin Chen, Jonas Kohler, and Michael Cohen.
2025. ``Storybooth: Training-free Multi-Subject Consistency for Improved
Visual Storytelling.'' arXiv:2504.05800.
https://arxiv.org/abs/2504.05800.

Tewel, Yoad, Omri Kaduri, Rinon Gal, Yoni Kasten, Lior Wolf, Gal
Chechik, and Yuval Atzmon. 2024. ``Training-Free Consistent
Text-to-Image Generation.'' arXiv:2402.03286.
https://arxiv.org/abs/2402.03286.

\begin{enumerate}
\def\labelenumi{\arabic{enumi}.}
\tightlist
\item
  \begin{enumerate}
  \def\labelenumii{\arabic{enumii}.}
  \setcounter{enumii}{2021}
  \tightlist
  \item
    Instagram post: ``a plate of food'' --- ``an exquisite corpse
    collaboration with @mrscottroberts through \#MidJourney.'' August
    16, 2022. https://www.instagram.com/p/ChWLLqvuOXe/.
  \end{enumerate}
\end{enumerate}

\end{document}
